\section{A complete source family separating eager and lazy replay}
\label{sec:separation}


\begin{theorem}[A native AHT star family]
\label{thm:star}
For integers $H\ge1$ and even $L\ge2$, set $N=(H+1)L$ and supply the inclusive
interval pairings
\[
 g_j:[0,L-1]\longrightarrow[jL,(j+1)L-1],\qquad 1\le j\le H.
\]
Assign weights $w(x)=x\bmod2$ on $[0,L)$ and $w(x)=1$ on $[L,N)$.
The maintained forward AHT scheduler produces exactly $H$ truncations at
$HL,(H-1)L,\ldots,L$, followed by one static contraction. Each truncated
suffix meets a single weight run, while every receiving interval contains
$L$ alternating runs. Its final histogram is
\[
 \{(H,L/2),(H+1,L/2)\},
\]
where pairs denote (weight, orbit multiplicity).

The previous eager native weight evaluator performs $\Omega(HL)$ run visits
or receiving-value updates on this family. Lazy evaluation uses
$O(L+H\log(L+2))$ weight-tree structural work after canonicalization, including
final emission. In a record-cost model charging one unit per dictionary
operation, generic source canonicalization takes $O(L\log(L+2))$ work,
and the unchanged unweighted trace checker takes $O(H^2)$ pairing-row
checks. In that model complete certificate replay costs
\[
 O\bigl(L\log(L+2)+H\log(L+2)+H^2\bigr).
\]
It has an asymptotic record-work separation from the eager method, for
example $\Omega(H^3)$ eager work versus $O(H^2\log(H+2))$ lazy work and
shared overhead when $L=2H^2$. The same upper bound has a deterministic
comparison-model realization using sorted endpoint preparation; on this
family the final histogram has only two keys. This final whole-verifier
bound counts map operations in the delivered Python implementation;
it does not assert a worst-case constant cost for its hash lookups or
for arithmetic on binary integers. The preceding weight-tree bound is
deterministic without that map-cost assumption.
\end{theorem}
\begin{proof}
All blocks $[jL,(j+1)L)$ are covered by pairing domains or ranges, so no
initial static gap exists. Only $g_1$ is periodic: its period is $L$ and its
domain and range are adjacent. For $j>1$ the distance $jL$ exceeds the domain
width, so two periodic pairings can never be merged. The unique pairing
whose range ends at $N-1$ is $g_H$, and every other pairing range ends below
$HL$. None is transmitted by $g_H$. The suffix exposed for truncation is
therefore exactly $[HL,(H+1)L)$. The carrier period is $HL$, its nominal
receiving interval is $[0,HL)$, and its transported addition is one on
$[0,L)$ and zero on $[L,HL)$. Thus only the first block changes, and
$g_H$ disappears.
The same argument applies inductively with $H-1$ in place of $H$. At the end
no pairing remains and the whole residual prefix contracts.

Every removed block still has constant value one: preceding folds changed
only the first block. After $t$ folds the first block has value
$t+(x\bmod2)$, so it retains all $L$ alternating runs. Initially its final
value one coalesces with the long suffix; after the first fold there is at
most one additional run. In all cases a removed suffix intersects one run.
Every orbit is $\{r,r+L,\ldots,r+HL\}$ for a unique $0\le r<L$, proving
the histogram formula.

The former flat evaluator scans the retained prefix and visits its $L$
alternating entries at every fold; the local-overlay producer also touches
the first block's $L$ entries to add one. This gives the stated
$\Omega(HL)$ cost.
In the lazy tree, one constant vector is added to the root of the first-block
subtree at each fold, with logarithmic splits and joins and no scan of that
subtree. The final contraction scans its $L$ leaves once. Initial source
and certificate-weight checks process $O(L)$ records. The unweighted
checker visits exactly $H+(H-1)+\cdots+1=H(H+1)/2$ current pairing rows
in its truncation precondition loops. Its only gap sort occurs after the
pairing list becomes empty; the carrier is always the last list element,
so its deletion causes no row shifts. Thus no extra $\log H$ factor is
needed for this checker. Generation of the trace is a separate phase and
is not included in this replay estimate. Add initialization and the
unchanged input/trace costs to obtain the stated bound.
\end{proof}

For $H>1$ the actual transport counters on this family are
\[
 \begin{aligned}
 \text{source runs}&=H,&
 \text{receiving runs}&=HL+H-2,\\
 \text{addition intervals}&=2H-1,&
 \text{emitted runs}&=L.
 \end{aligned}
\]
The maximum canonical run count is $L+1$. For $H=1$, the receiving and
maximum counts are both $L$. These receiving counts are logical state
sizes read from cached subtree metadata, not counts of visited leaves.

This family was checked against the maintained producer and independent
weighted verifier for $(H,L)=(1,2),(3,8),(8,64),(16,256)$: the complete operation
lists, cut positions, and final histograms agree with the theorem. These
finite executions check compatibility with the current scheduler; the proof
above establishes the arbitrary-parameter family. The examples are synthetic
interval-pairing inputs, not claimed difficult knot diagrams.
